\section{Conclusion}\label{sec:conclusion}

This paper addresses two fundamental challenges in graph-based traffic forecasting: the oversmoothing caused by dense physical graphs, and the temporal heterogeneity of traffic dependencies that a single static graph cannot capture.

Our key findings are:

\begin{enumerate}
    \item Dense physical road-network graphs can severely degrade GCN/T-GCN forecasting performance through excessive spatial aggregation. On the Los-loop dataset, removing the physical graph entirely improves RMSE by 32.8\%.

    \item Traffic dependencies are temporally heterogeneous: different temporal lags exhibit distinct functional dependency structures with low cross-lag overlap.

    \item Multi-lag DAGMA explicitly discovers these lag-specific dependencies by constructing temporal input blockings $Z = [x(t{-}L), \ldots, x(t)]$ and extracting separate functional graphs for each lag.

    \item T-GCN-MultiGSL-Mix exploits these lag-specific graphs through per-node, per-timestep learned mixing, achieving a mean RMSE improvement of 14.9\% over a no-graph baseline on Los-loop across five random seeds at 5-minute sampling, and 27.4\% at 15-minute sampling on the same network.

    \item The improvement is robust: it holds in all five seeds at both sampling intervals, survives a parameter-matched control, and persists across all four prediction horizons.

    \item The effect is dataset-dependent: on SZ-Taxi, improvements are marginal ($\leq0.3\%$, with $-0.02\%$ at PH=4), indicating that the multi-lag approach is most effective when traffic networks exhibit strong lag-specific propagation patterns. Temporal resolution is ruled out as the explanation, as Los-loop improves strongly at both 5-minute and 15-minute sampling.
\end{enumerate}

Future work should explore learned lag window selection, sliding-window DAGMA for time-varying graph updates, scalability to larger networks, and integration with other graph-based forecasting architectures.
